% --- Avant-propos -----------------------------------------------------------
\chapter*{Avant-propos}
\markboth{Avant-propos}{Avant-propos}
\addcontentsline{toc}{chapter}{Avant-propos}

Ce livre est le troisième de la collection Boky Matematika, après Terminale S
et Terminale OSE. Il en reprend les principes sans en changer un seul : un
même texte pour l'élève et pour le professeur, un cours qui ne mobilise
jamais une notion avant de l'avoir enseignée, et une banque d'exercices
exclusivement tirée de sujets réellement posés, chacun avec sa provenance.

Il s'en écarte sur un point, et ce point commande tout le reste : le
programme de la série littéraire tient en \textbf{deux heures par semaine}.
C'est moins de la moitié du temps dont dispose la série S. On ne peut donc
pas faire tenir dans ce livre ce qu'un livre de série S contient — et il
serait absurde d'essayer. Ce que l'on peut faire, en revanche, c'est
travailler chaque notion beaucoup plus souvent, sous des formes plus variées,
jusqu'à ce qu'elle devienne un réflexe. C'est le parti pris de ce volume :
peu de notions, beaucoup d'exercices, et des formats différents — questions à
choix multiples, vrai ou faux, automatismes, lectures graphiques, situations
concrètes, annales, devoirs surveillés — pour que la même idée soit rencontrée
sous plusieurs angles avant d'être vue en épreuve.

\section*{Ce que dit le programme officiel, et ce que disent les sujets}

Le Programme d'Études 2026 décrit pour la série littéraire trois blocs de
contenus : l'analyse (trente heures), l'algèbre (douze heures), le traitement
des données et les probabilités (vingt-deux heures). L'analyse y va plus loin
qu'on ne l'attend d'un horaire de deux heures : elle comprend le logarithme
népérien, l'exponentielle, la primitive et l'intégrale.

Les suites numériques, elles, n'apparaissent pas dans ce tableau des
contenus. Elles figurent en revanche dans la Répartition Annuelle, qui leur
consacre la quatrième période entière, dans le programme officiel de 2019,
qui leur accorde six semaines — et surtout dans les sujets : le premier
exercice du baccalauréat porte sur les suites à chacune des sessions
examinées pour ce livre, de 1999 à 2026, sans une seule exception. Il n'était
pas envisageable de les laisser de côté. Elles forment donc ici une
\textbf{partie à part entière}, la troisième, avec sa couleur propre : c'est
la façon la plus honnête de signaler qu'elles viennent de la Répartition
Annuelle et des sujets, et non du tableau des contenus du PE.

Symétriquement, le dénombrement — arrangements, permutations, combinaisons —
n'est pas listé comme un contenu séparé du PE, mais aucun exercice de
probabilité de la série n'est traitable sans lui, tirage simultané et tirages
successifs étant explicitement au programme. Il ouvre donc la quatrième
partie.

\section*{Un mot à l'élève de terminale}

Vous êtes peut-être entré en série littéraire parce que les mathématiques ne
vous attiraient pas. Le programme qui vous attend n'a rien d'écrasant : trois
grands objets, une fonction que l'on étudie, une suite dont on cherche le
terme général, un tableau de chiffres dont on tire une prévision. En deux
heures par semaine, tout tient — à condition de ne rien laisser passer.

Vous remarquerez que chaque chapitre commence par une page de révision qui ne
parle jamais de la notion nouvelle. C'est voulu : on ne vous demandera jamais
de deviner ce que vous n'avez pas encore appris. Vous remarquerez aussi que
les exercices ne se ressemblent pas d'une page à l'autre. Un QCM ne se
travaille pas comme une annale, et une lecture graphique ne se travaille pas
comme un calcul de dérivée. C'est en changeant de format que l'on découvre ce
que l'on croyait savoir sans le savoir.

Le barème compte. L'épreuve dure deux heures quinze, les deux exercices
valent cinq points chacun et le problème dix. Cela veut dire qu'un exercice
vaut, à lui seul, un quart de la note. Les élèves qui échouent ne sont
presque jamais ceux qui n'ont pas su faire le problème : ce sont ceux qui ont
laissé filer les deux exercices.

\section*{Un mot au professeur}

Le livre est rangé par domaine, pas dans l'ordre où l'on enseigne. L'annexe
de progression, réservée à l'édition du professeur, propose la conversion en
cinq périodes, avec une colonne vide pour vos dates. Elle ne porte aucune
date : ce livre est fait pour resservir d'une année sur l'autre.

L'ordre des parties mérite une justification. La partie sur les suites vient
après l'analyse, alors que la Répartition Annuelle place les suites en
quatrième période et le logarithme nulle part. La raison est la règle
d'antériorité, tenue ici sans exception : les sujets écrivent régulièrement
$V_{n} = e^{2-n}$, $V_{n} = e^{3n-1}$ ou $W_{n} = \ln(U_{n})$, et l'on ne
peut pas poser ces exercices à un élève qui n'a pas encore rencontré le
logarithme et l'exponentielle. En classe, rien n'empêche de traiter les
suites plus tôt : il suffit de s'en tenir aux exercices qui ne font pas appel
à ces deux fonctions, et le livre signale lesquels.

Les deux options de la série, A1 et A2, passent la même épreuve, mais A2 y
répond à des questions supplémentaires — presque toujours une primitive et un
calcul d'aire — et sur un barème différent, de coefficient 3 contre 1. Les
questions qui leur sont réservées sont signalées dans tout le livre par une
marque explicite. Elles ne sont jamais isolées en fin de chapitre : un élève
de A1 doit pouvoir sauter la question sans perdre le fil, et un élève de A2
doit la trouver à sa place logique.

Aucun sujet de concours ENI ou ESPA ne figure ici : ces concours recrutent
sur un profil scientifique sans rapport avec la série littéraire. La banque
d'annales s'appuie sur les sessions malgaches de la série~A, de 1999 à 2026,
et sur les baccalauréats blancs conservés.

Les démonstrations sont écrites. Cela ne veut pas dire qu'elles doivent
toutes être faites au tableau — deux heures par semaine ne le permettraient
pas. Elles sont là pour que l'élève qui veut comprendre pourquoi un résultat
est vrai trouve la réponse dans son livre plutôt que nulle part.

\vspace{1.2em}
\noindent\textit{Antananarivo}
